\chapter{Átomos y Tipos de Datos Primitivos}
\label{chap:atoms}

{\huge$\bigstar$} Principiante

Todo valor en LisPy es, en última instancia, un átomo o una celda cons. Este
capítulo cubre los tipos atómicos fundamentales.

\section{Números Enteros}

Los enteros en LisPy tienen precisión arbitraria (bigints):

\begin{lispcode}
42
-7
123456789012345678901234567890
\end{lispcode}

No hay desbordamiento: un entero puede tener tantos dígitos como la memoria
lo permita.

\section{Números Racionales}

Los racionales se escriben como fracción:

\begin{lispcode}
1/3
22/7
\end{lispcode}

Las operaciones aritméticas entre racionales producen resultados exactos:

\begin{lispcode}
(+ 1/3 1/7 1/11)  ;; → 137/231
\end{lispcode}

\section{Números Flotantes}

Los flotantes siguen el estándar IEEE 754 de doble precisión:

\begin{lispcode}
3.14159
1.0e10
-0.5
\end{lispcode}

\section{Números Complejos}

Los números complejos se construyen con la notación \texttt{\#C}:

\begin{lispcode}
#C(3 4)       ;; 3 + 4i
#C(1.0 -2.5)  ;; 1.0 - 2.5i
\end{lispcode}

\section{Caracteres}

Los caracteres se escriben con la sintaxis \texttt{\#\textbackslash}:

\begin{lispcode}
#\A          ;; la letra A mayúscula
#\Space      ;; el carácter de espacio
#\Newline    ;; el carácter de nueva línea
#\λ          ;; lambda griega (Unicode)
\end{lispcode}

\section{Strings}

Las cadenas se delimitan con comillas dobles:

\begin{lispcode}
"Hello"
"Línea 1~%Línea 2"  ;; ~% inserta un salto de línea
\end{lispcode}

\section{Símbolos}

Los símbolos son identificadores con nombre. Se escriben con comilla simple
para evitar su evaluación:

\begin{lispcode}
'foo                      ;; el símbolo FOO
'|nombre con espacios|    ;; símbolo con espacios internos
:keyword                  ;; keyword (símbolo en el paquete KEYWORD)
\end{lispcode}

\section{Booleanos}

\begin{lispcode}
t    ;; verdadero
nil  ;; falso (también la lista vacía)
\end{lispcode}

En Common Lisp, \texttt{nil} es el único valor falso; cualquier otro valor es
verdadero en contexto booleano.

\section{Conversión entre Tipos}

\begin{lispcode}
(float 42)           ;; → 42.0
(coerce 3/2 'float)  ;; → 1.5
(numerator 3/2)      ;; → 3
(denominator 3/2)    ;; → 2
\end{lispcode}

\section{Ejemplos}

\begin{examplebox}[title=Precisión arbitraria]
\begin{lispcode}
;; Fibonacci(100)
(defun fib (n)
  (if (< n 2) n
    (+ (fib (- n 1)) (fib (- n 2)))))
;; Nota: usar versión iterativa para valores grandes
\end{lispcode}
\end{examplebox}

\begin{examplebox}[title=Aritmética exacta con racionales]
\begin{lispcode}
(+ 1/3 1/7 1/11)  ;; fracción exacta: 137/231
(* 1/3 3)         ;; → 1
\end{lispcode}
\end{examplebox}

\begin{examplebox}[title=Números complejos]
\begin{lispcode}
(* #C(3 4) #C(1 -2))  ;; → #C(11 -2)
(abs #C(3 4))         ;; → 5.0 (magnitud)
(phase #C(1 1))       ;; → 0.785398... (fase en radianes)
\end{lispcode}
\end{examplebox}

